\section{Architect-Identified Algorithmic Strategies Across Domains}
\label{app:architect-strategies}

This appendix summarizes the blueprints produced by the Architect for the three
NaturalPlan domains.  In our Architect--Builder--Runner setup, the Architect is
not asked only to describe the task.  It must specify an executable planning
design: the representation of the problem, the functions that the Builder
should implement, the algorithmic strategy for finding a plan, and the logic
that connects parsing, solving, validation, and output generation.  The
following sections make these domain-specific blueprints explicit.

\subsection{Trip Planning}

The Architect treats trip planning as a constrained path-finding problem over a
flight graph.  Cities are graph nodes, direct flights are directed edges, and a
complete itinerary is an ordered path that visits every required city exactly
once.  The temporal structure is represented by inclusive visit intervals.  If
the traveler flies from one city to the next on day \(d\), then day \(d\) is
shared by the two consecutive city intervals.  Therefore, an itinerary with
\(n\) cities must satisfy the total-day identity
\(\sum_i \mathrm{duration}_i - (n-1)\).

This blueprint specifies:
\begin{itemize}
    \item \textbf{State representation:} a problem instance contains the city
    list, required stay duration for each city, total trip length, directed
    flight edges, and fixed event windows.  A partial or complete plan is a
    sequence of visit records \((city, start, end)\).  The state must preserve
    city coverage, valid flight connectivity, inclusive day intervals, and the
    requirement that every fixed event window lies inside the visit interval
    for its city.

    \item \textbf{Core functions:} the Builder is expected to implement
    parsing routines for cities, durations, total days, flight edges, and event
    constraints; candidate-generation logic for extending a partial route to
    feasible next cities; consistency checks for total-duration arithmetic; a
    route verifier that checks city coverage, flights, intervals, and events;
    and a formatter that converts the verified route into a day-by-day trip
    plan.

    \item \textbf{Solving strategy:} the Architect specifies deterministic
    constrained depth-first search over possible flight paths.  The search
    prunes a branch when the current city interval violates an event window,
    exceeds the trip horizon, or makes a remaining city impossible to schedule
    later.  For small underconstrained cases, the blueprint permits a
    deterministic permutation-based fallback over city orders.

    \item \textbf{Integration logic:} the system first parses the natural
    language prompt into a structured graph-and-interval problem.  It then
    checks the total-day arithmetic induced by shared flight days, searches for
    a feasible city ordering, verifies the resulting itinerary against all
    parsed constraints, and returns the plan only after this verification
    succeeds.
\end{itemize}

\subsection{Meeting Planning}

The Architect treats meeting planning as a time-window routing problem.  The
agent starts at a specified location and time, then chooses an ordered sequence
of friends to meet.  Each transition must account for travel time, possible
waiting, the friend's availability window, and the minimum required meeting
duration.  Unlike trip planning, the goal is not necessarily to visit every
friend; the objective is to maximize the number of feasible meetings.

This blueprint specifies:
\begin{itemize}
    \item \textbf{State representation:} a planning state consists of the
    current location, current clock time, and set of already-met friends.  Each
    friend is represented by a name, location, availability interval, meeting
    duration, and stable input index.  The travel-time matrix defines the cost
    of moving between locations.  A candidate meeting is valid only if the
    agent can arrive, wait if needed, and complete the meeting before the
    friend's availability window closes.

    \item \textbf{Core functions:} the Builder is expected to implement time
    parsing and formatting, availability-window parsing, construction of the
    structured meeting problem, travel-time lookup, candidate generation over
    unmet friends, recursive schedule search, tie-breaking logic, schedule
    verification, and natural-language plan formatting.  Each candidate action
    records the departure location, travel time, arrival time, meeting start,
    and meeting end.

    \item \textbf{Solving strategy:} the Architect specifies exact dynamic
    search with memoization over subset-style routing states.  From each state,
    the solver evaluates all friends who have not yet been met and recursively
    computes the best continuation.  The primary objective is to maximize the
    number of meetings.  Ties are resolved deterministically, favoring earlier
    completion, less waiting, and stable input order.

    \item \textbf{Integration logic:} the system parses the start state,
    friend constraints, and travel matrix; computes the best feasible meeting
    sequence; replays the sequence to verify travel times, arrival times,
    waiting, availability windows, meeting durations, and duplicate-meeting
    violations; and finally emits the verified route as an ordered list of
    plan steps.
\end{itemize}

\subsection{Calendar Scheduling}

The Architect treats calendar scheduling as finite interval enumeration over an
ordered work-calendar grid.  The problem is to choose one meeting slot that
fits within the allowed days and work hours while avoiding all participant
conflicts and respecting any stated preferences.  Because the candidate days,
work hours, and meeting duration are bounded, the Architect chooses direct
constraint enumeration rather than a more complex search procedure.

This blueprint specifies:
\begin{itemize}
    \item \textbf{State representation:} a calendar instance contains the
    participants, ordered candidate days, work-hour start and end times,
    required meeting duration, per-participant busy intervals, blocked days,
    lower time bounds, upper time bounds, and an earliest-availability
    preference.  A candidate plan is a tuple \((day, start, end)\), where the
    interval length must equal the requested duration.

    \item \textbf{Core functions:} the Builder is expected to implement clock
    parsing, duration parsing, day extraction, participant extraction,
    busy-interval parsing, interval merging, preference parsing, candidate-slot
    generation, slot-validity checking, result verification, and output
    formatting.  The candidate generator scans possible starts on a
    deterministic 30-minute grid.

    \item \textbf{Solving strategy:} the Architect specifies chronological
    constraint enumeration.  For each allowed day, the solver enumerates
    candidate start times and rejects a slot if it falls outside work hours,
    overlaps any participant's busy interval, occurs on a blocked day, starts
    before a parsed lower bound, or ends after a parsed upper bound.  In the
    default case, the first feasible slot is returned; when the parsed
    preferences call for a later free block, the solver deterministically
    selects the beginning of the last available block.

    \item \textbf{Integration logic:} the system parses the prompt into a
    structured calendar problem, normalizes busy intervals by merging overlaps,
    enumerates valid slots in day and time order, constructs the proposed
    meeting time, verifies its duration and participant-level constraints, and
    returns the final slot only if all checks pass.
\end{itemize}

Overall, the Architect's blueprints identify three distinct algorithmic
structures for NaturalPlan: constrained graph search for trip planning,
memoized time-window routing for meeting planning, and interval-based slot
enumeration for calendar scheduling.  These blueprints make clear what the
Architect actually specifies: the planning state, the required parsing and
validation functions, the candidate-generation mechanism, the solving
algorithm, and the final integration path from natural-language prompt to
verified plan.
